% ===== §22 =====
\section{High-Density Sedimentation and Accelerated Time}
\label{sec:sedimentation}

\noindent
The third pillar concerns the temporality of intimate culture: the speed at which it emerges. The previous pillars established what is woven, in the vacancy left by the absent big Other; this one establishes how fast the weaving sediments into durable form, and shows that the answer is peculiar to the dyad and has a consequence no other cultural formation affords---that the moment of cultural emergence is brought within the horizon of lived experience, so that the couple are not only the bearers of their culture but the witnesses of its birth.

\subsection*{The density of shared living}

The rate at which a culture sediments is a function of the density of the shared living from which it precipitates, and the intimate dyad is the densest medium of shared living there is. The density has several dimensions, and they compound. There is the density of \emph{time}: intimate partners share not scheduled intervals but the continuous duration of a shared life, mornings and nights and the unstructured hours between, a volume of shared time no other relation approaches. There is the density of \emph{body}: intimate relation is carried on in bodily proximity, in touch, in the sharing of space and sleep and the physical routines of a life, so that its medium is not only spoken but corporeal. There is the density of \emph{frequency}: the interactions are not occasional but constant, iterated many times a day, so that any form has abundant occasions to be repeated and confirmed. And there is the density of \emph{stakes}: what passes between intimate partners matters to them as almost nothing else does, so that each interaction carries an affective charge that fixes it in memory and makes it a candidate for repetition. These densities together make the intimate dyad a medium in which the precipitation of shared living into cultural form occurs faster than anywhere else, because every condition that favors sedimentation---time, repetition, bodily co-presence, affective stakes---is present at its maximum.

\subsection*{Emergence brought within the horizon of experience}

The consequence is the distinctive feature of this pillar: intimate culture emerges at a speed that brings the moment of emergence \emph{within the horizon of lived experience}. In macro-culture the emergence of a cultural form is never witnessed; it occurs across generations, below the threshold of any individual's experience, so that the form is always already there when any given person meets it, and its origin is a matter for historical reconstruction rather than memory. No one witnessed the birth of a national idiom or the first performance of a folk rite; these emerged too slowly and too diffusely for emergence to be an object of experience. The intimate dyad is different in kind. Because its culture sediments at the speed the density affords, the couple can often witness, and later recall, the very moment a form was born: the occasion on which a private word was coined and, in the same evening, accepted; the first time a gesture was made that would become a rite; the moment a shared glance was found to carry a meaning both had, without saying so, arrived at. The emergence of culture, which in every larger formation lies below the horizon of experience, rises in the dyad above it. This is why the intimate dyad is the privileged site for a phenomenology of cultural emergence: it is the one cultural formation in which emergence is not reconstructed but witnessed, the one place where one can watch a culture being born.

\subsection*{The moments of self-encoding}

These witnessed moments have a structure worth describing, for they are the phenomenological appearance of the exteriorization the first part defined, caught in the act. What happens, in such a moment, is that something that occurred once---a contingent event, an accident, a joke that landed, a mistake, a shared perception---is, by both parties at once, taken up and marked as \emph{significant}, as belonging to them, as a form to be kept and repeated rather than a happening to be let go. The moment has a double structure. First there is the contingent occurrence, which might have passed and been forgotten. Then there is the mutual, often silent, recognition---a look exchanged, a laugh that both know they will return to, a repetition offered by one and accepted by the other---in which the occurrence is lifted out of the flow of forgettable events and installed as a form. This second act is the moment of \emph{self-encoding}: the relation, in it, encodes a piece of its own generative history into a form it can re-enact, and does so through the simultaneous recognition of both parties. It is the Hegelian exteriorization made visible: the relation goes out of itself into a form, in the very instant of the mutual recognition that constitutes the form, and one can see it happen. That the encoding requires both parties at once---that a private form cannot be instituted by one alone but only by the two together, in a recognition each reads in the other---is the trace, at the level of the emergence-moment, of the vacancy of the big Other established in the first pillar: with no external authority to institute the form, the two must institute it themselves, together, and the moment of self-encoding is the moment they do.

\subsection*{Speed and fragility}

One consequence of this accelerated sedimentation must be marked, because it qualifies the pillar and connects it to the fifth. The same speed that lets intimate culture emerge within the horizon of experience also makes it, in a certain sense, thin in its foundations. A macro-culture's forms are sedimented over so long a time and through so many bearers that no single event and no single defection can unsettle them; they have the stability of the deeply laid. The dyad's forms are laid down fast, by two people, over a time short enough to remember, and what is deposited quickly can be disturbed quickly. The accelerated emergence that makes the dyad the privileged site of witnessed cultural birth is the same accelerated temporality that leaves its culture without the deep foundations time alone can give, and therefore exposed, as the fifth pillar will show, to a reversibility no macro-culture faces. The speed is a gift and a vulnerability at once: it brings emergence within experience, and it denies the culture the ballast of slow sedimentation. Having established what is woven, that it spans the registers, and the speed at which it sediments, the phenomenology turns to the register in which all of this is most intensely configured, before it turns, finally, to the fragility the speed has already foreshadowed.
